\cotefile{157-5}
\scanpage{421}
\mspage{333}

\[
\begin{aligned}
x&=\varinjlim_{i\in I}x_i,\qquad \text{les }x_i\text{ dans }\underline S(\pi)\\
&& I\ \unsure{\text{grd}|\text{gd}|\text{gros}}\text{ devant }\pi.
\end{aligned}
\]

Alors le foncteur composé $i\mapsto x_i$
\diagram{\begin{tikzcd}[ampersand replacement=\&]
 I \arrow[r] \& \underline S(\pi)_{/x} \arrow[r,hook] \& \unsure{M|\mathcal M}_{\pi/x}
\end{tikzcd}}
est cofinal\inserted{(cf. 4.4.2 \unsure{pag.|page|p.} 79)}, \deleted{et comme} \inserted{donc} $\underline S(\pi)_{/x}$ \unsure{$-\!-$|$\longrightarrow$} $M_{\pi/x}$ \deleted{est \unsure{pleint|pleint\textsuperscript{t}|plein\textsuperscript{t}} fidèle, on voit qu'il} satisfait la condition F1 (puisque le composé la satisfait), et comme $\underline S(\pi)_{/x}\longrightarrow M_{\pi/x}$ est pl. fid, cela implique déjà que ce foncteur est cofinal. Dans \deleted{\unsure{par}} \unsure{cas}, on en conclut $x\in\mathrm{Ob}\,\underline S$, qed.

\deleted{\unsure{D}}Supposons maintenant que $S$ satisfasse \unsure{a')|d')} (au lieu de b)), et prouvons encore l'équivalence \unsure{$(*)$|$(**)$}. Pour ceci, écrivons
\[
 x=\varinjlim_{j\in\unsure{J|I}}x_j\qquad
 \left|\begin{aligned}
 &\text{les }x_j\in\mathrm{Ob}\,M_\pi\\
 &J\ \text{\uline{ens. ordonné}}\quad\unsure{*}\\
 &(x_{\unsure{i|j}})_{\unsure{i|j}\in\unsure{J|I}}\ \deleted{\text{est}}\ \text{est }\unsure{\pi_0|\text{bon}} \text{-}\pi\text{-adapté}\\[-2pt]
 &\deleted{\unsure{=|-}\,\unsure{\text{grd}|\text{gd}}\text{ devant }\pi}
 \end{aligned}\right.
\]
% slanted
\margin{\uline{NB} Ça existe, \unsure{grâce|après|appliq.} \unsure{à|?} 4.13.14. $\square$}
ce qui est possible puisque \unsure{$\pi$|il} est \unsure{utile}, pour $M$. \unsure{$]$} Définissons $J_{\underline S}$ comme dans l'énoncé de a), donc comme image inverse de $\underline S(\pi)_{/x}$ $\unsure{\subset|\hookrightarrow}\,M_{\pi/x}$ par $I\unsure{\longrightarrow|-}M_\pi$. Par hyp. \unsure{d')|a')}, on a
\[
x\unsure{\mathord{+}|\in}\mathrm{Ob}\,\underline S\quad\Longleftrightarrow\quad I_{\underline S}\text{ est cofinal dans }I
\tag*{$(**')$}
\]

En utilisant $(**')$, prouvons \inserted{\unsure{encore|alors}} $(*)$\inserted{(p. \deleted{\unsure{0}}\unsure{331})}. \deleted{\unsure{P}} En fait, on utilisera \unsure{seulement} qu'on\unsure{.| a}

% archive sheet number: 420

\scanpage{422}
\mspage{334}

\deleted{\unsure{prouvons|prouver|prouve} $\Longrightarrow$, supposons que}

l'isomorphisme $\square$, avec $I$ \unsure{grd|gd|gros} devant $\pi$, et qu'on a l'équivalence $(**')$, sans avoir à utiliser que $I$ est \unsure{ordonnée|ordonné}. Donc la démonstration vaut aussi dans le cas c).

Prouvons $\Longrightarrow$ dans $(*)$, i.e. supposons $x\in\mathrm{Ob}\,\underline S$, prouvons que $x$ est dans l'image essentielle de $\mathrm{Ind}_\pi(\underline S(\pi))$. En vertu de \unsure{$(**')$|$(**)$}, $I_{\underline S}$ est cofinal dans $I$, et comme $I_{\underline S}\unsure{\longrightarrow|-}I$ est une inclusion pl. \unsure{fidle|fidèle|fid.}, il s'ensuit que $I_{\underline S}$ est également grd devant $\pi$, d'autre part on aura par cofinalité
\[
x=\varinjlim_{i\in I_{\underline S}}x_i,
\]
où cette fois les $x_i$ sont dans $\underline S(\pi)=\underline S\cap\unsure{M|\mathcal M}_\pi$, d'où résulte que $x$ est \unsure{que|dans} l'image ess. de $\mathrm{Ind}_\pi(\underline S(\pi))$.

Prouvons $\Longleftarrow$ dans \unsure{$(*)$}, i.e. supposons que \deleted{l'on ait $x$} l'on ait \deleted{\unsure{$x=$}}
\[
\begin{aligned}
&\deleted{x\unsure{=x|\!-\!x}}\\[-2pt]
&x=\varinjlim_{j\in J}x_j\qquad
\left|\begin{aligned}&\text{les }\unsure{x_j|y_j}\text{ dans }\underline S(\pi)\\
&J\text{ grd devant }\pi.\end{aligned}\right.
\end{aligned}
\]

% archive sheet number: 421

\scanpage{423}
\mspage{335}

En invoquant à nouveau 4.4.\unsure{2|1}, on trouve que $J\longrightarrow\unsure{M|\mathcal M}_{\pi/x}$ est cofinal, et comme ce foncteur se factorise par le foncteur d'inclusion \unsure{pl fidèle|pl. fidèle} $\underline S(\pi)_{/x}\hookrightarrow M_{\pi/x}$, il s'ensuit \deleted{que} comme tantôt que ce dernier est cofinal. Nous allons \deleted{\unsure{en conclure}} \inserted{admettre} pour un moment que ceci implique \unsure{que|pour} l'image inverse \deleted{$\unsure{J}$}$I_{\underline S}$ de cette sous-catégorie str. pleine, par le foncteur également cofinal $I\longrightarrow M_{\pi/x}$\unsure{;|,} est cofinal dans $I$. Ceci étant, l'équivalence $(**')$ (p. 333) implique $x\in\mathrm{Ob}\,\underline S$, ce qu'on voulait.

Donc il reste à prouver \deleted{\unsure{que le}} que si
\diagram{\begin{tikzcd}[ampersand replacement=\&]
 J'=\underline S(\pi)_{/x}\arrow[r,hook] \& M_{\pi/x}\unsure{=|\subset}J
\end{tikzcd}}
est cofinal, alors $I'=I_{\unsure{S|\underline S}}\longrightarrow I$ l'est aussi. Il semble qu'il faille \inserted{bel et bien} \unsure{(à nouveau|à nouveau} \deleted{\unsure{la}} invoquer les conditions de stabilité \deleted{\unsure{de}} $(\pi_0,\pi)$ \unsure{faite|faites} dans \unsure{d')|a')} (ce qui n'avait pas été nécessaire dans le cas c)), \unsure{et}

% archive sheet number: 422

\scanpage{424}
\mspage{336}

le lemme\inserted{-clef} (\unsure{un peu}) vache 4.13.17, où on prendra (comme dans la démonstration \unsure{décrite|écrite} avant l'énoncé du dit lemme)
$\alpha=\alpha_\pi\overset{\text{déf}}{=}$ plus petit ordinal $\alpha$ tel que $\mathrm{card}\,\alpha=\pi$. Donc $I_\alpha$ est de cardinal $\pi$, et \unsure{grd|gd|gros} devant \unsure{tt|tout} cardinal $\pi'<\pi$, et en particulier devant $\pi'_0$. L'hypothèse $2^\circ)$ de 4.13.17 est vérifiée sur $I$, par définition d'un syst. inductif $\pi_0$-$\pi$-adapté, et pour $J'$, grâce à l'hyp. $\mathrm L(\pi_0,\pi)$ sur $\underline S(\pi)$. De $=$, l'hyp. $3^\circ)$ \inserted{($I$ grd devant \deleted{\unsure{tout}} cardinal $<\pi$)} est satisfaite par \unsure{définition} des systèmes $\pi_0$-$\pi$-adaptés -- $I$ \deleted{\unsure{sont|soit}} $=$ grd devant $\pi$. L'hypothèse $1^\circ)$ a déjà été vue.

On a donc prouvé (par voie circulaire)
\diagram{\begin{tikzcd}[ampersand replacement=\&,column sep=small,row sep=small]
 \& \& \& \deleted{a}\unsure{c'|a'} \& {}\\
 a \arrow[r,Rightarrow] \& {(\deleted{a}\unsure{b'|a'}} \arrow[r,Leftrightarrow] \& \deleted{a}\unsure{c|a}\arrow[ur,Leftrightarrow]\arrow[dr,Rightarrow] \& \& {),}\\
 \& \& \& b \& {}
\end{tikzcd}}
en ayant à invoquer le lemme-clef ci-dessus pour les \uline{deux} implications
\[
 \deleted{a}b'\Longrightarrow\deleted{a}c\qquad\text{et}\qquad\deleted{a}\unsure{c'|a'}\Longrightarrow\deleted{a}b'
\]
(donc pour deux implications \uline{en sens inverse}, chose \unsure{étrange} \unsure{\ldots}).

% archive sheet number: 423

\scanpage{425}
\mspage{337}

Il reste seulement encore à prouver l'implication
\[
\deleted{a}\unsure{b'|a'}\Longrightarrow a,
\]
ce qui est essentiellement trivial.

Cela achève la démonstration du critère \deleted{\unsure{de}}d'accessibilité \inserted{4.13.15} d'une sous-catégorie pleine $\underline S$ de $\unsure{M|\mathcal M}$.

\deleted{\uline{Corollaire} 4.13.18. Soient $M$ une catégorie \inserted{accessible}, \deleted{\unsure{un cardinal} tel \unsure{$\pi_0$}\inserted{\unsure{$\pi_0$-accessible}} que $M$ \unsure{satisfasse|satisfait} $\mathrm L_{\pi_0}$ ($M$ stable \unsure{$\pi$} \unsure{(un cardinal} \unsure{est} utile pour $M$, par $\varinjlim$ \unsure{grdes|grds|strictes} \unsure{devant|sur} $\pi_0$)}, $\pi$ un cardinal utile pour $M$, $\underline S$ une sous-catégorie str. pleine de $M$, $\underline S(\pi)=\underline S\cap M_\pi$. Soit $\alpha_0$ un ordinal limite, tel que $\mathrm{card}\,\alpha_0\leq\pi$, \deleted{et tel que $I_{\alpha_0}\unsure{:=|\!=}\unsure{\alpha<\alpha_0|[\alpha<\alpha_0]}$ soit \unsure{grd} devant $\pi_0$.} \deleted{et désignons} et tel \deleted{appelons \unsure{$M$} $\mathrm L(\unsure{\alpha|\vec\alpha})$ l'hypothèse : $\underline S(\pi)$ est} stable dans $M_\pi$ par $\varinjlim$ de type $I_{\alpha_0}$. (NB $M_{\unsure{\pi}}$ est stable par ces limites, et $M_\pi$ est stable dans $M$ pour ces limites). \deleted{\unsure{Donc} que $\underline S$}}

On va donner une variante un peu \unsure{décantée|décanté|différente} de 4.13.15, qui (avec des hyp. \unsure{un peu} plus faibles) capte l'essentiel de sa substance. L'énoncé 4.13.15 en est une conséquence \unsure{presque} immédiate :

% archive sheet number: 424

\scanpage{426}
\mspage{338}

\uline{Corollaire} 4.13.18. \inserted{(de 4.13.15)} Soient $\unsure{M|\mathcal M}$ une catégorie accessible, $\pi$ un cardinal utile pour $M$, $\alpha_0$ un \uline{ordinal} \uline{limite} tel que l'on ait $\mathrm{Card}\,\alpha_0\leq\pi$ \unsure{;} Soit $I_{\alpha_0}=\{\alpha\mid\alpha<\alpha_0\}$, et supposons que $M$ soit stable par $\varinjlim$ de type $I_{\alpha_0}$.

\figure{\begin{tikzpicture}[x=1em,y=1em]\draw (-.9,0)--(0,0)--(0,1.2)--(1.25,1.2);\draw (.11,0)--(.11,1.07);\end{tikzpicture}}

Soit donc $\underline S$ une sous-catégorie str. pleine de $M$, $\underline S(\pi)=\underline S\cap M_\pi$. Désignons par $\mathrm L(\alpha_0)$ la condition : $\underline S(\pi)$ est stable \deleted{par} \inserted{dans} $M_\pi$\inserted{(ou ce qui revient au $=$, dans \uline{$M$})} par $\varinjlim$ de type $I_{\alpha_0}$. \deleted{\unsure{Considérons}} \inserted{Alors} les conditions suivantes sont équivalentes

a) $\underline S$ est $\pi$-accessible dans $M$, et \deleted{\unsure{satisf}} est stable dans $M$ par $\varinjlim$ de type $I_{\alpha_0}$.

a') $\underline S$ est $\pi$-accessible et \inserted{$\underline S(\pi)$} satisfait $\mathrm L(\alpha_0)$.

\deleted{\unsure{a}}\unsure{b)} $\underline S(\pi)$ satisfait $\mathrm L(\alpha_0)$, et pour tout $x\in\mathrm{Ob}\,M$, \deleted{et toute \unsure{représentation} \inserted{\unsure{représentation} de $\pi$} \unsure{inductif} \unsure{pour} $x\simeq\varinjlim_{i\in I}x_{\unsure{j|i}}$, \unsure{avec} les $I$ \unsure{grd devant|grd et} $\pi$ et $x_i\in M_\pi$} on a ceci : $x\in\mathrm{Ob}\,\underline S$ \unsure{sss|ssi} la sous-catégorie str. pleine $\underline S(\pi)_{/x}$ \deleted{est} de $M(\pi)_{/x}$ est cofinale dans $M(\pi)_{/x}$.

\deleted{\unsure{b}}\unsure{c)} $\underline S(\pi)$ satisfait $\mathrm L(\alpha_0)$, de plus pour \deleted{\unsure{tt}} $x\in\mathrm{Ob}\,M$, et \deleted{toute} représentation
\[
 x\simeq\varinjlim_{\unsure{i\in I}}x_i\qquad
 \left|\begin{aligned}
 &\deleted{\unsure{x_i\in\mathrm{Ob}\,M_\pi}}\\
 &\deleted{\unsure{I\text{ grd devant }\pi}},\end{aligned}\right.
\]

% archive sheet number: 425

\scanpage{427}
\mspage{339}

qui est ``$\alpha_0$-$\pi$-adaptée'' i.e. telle que \deleted{\unsure{2) les $x_i\in$}} \deleted{1) les $x_i\in\mathrm{Ob}\,\unsure{M|\mathcal M}_\pi$ 2)} $1^\circ)$ $I$ \inserted{équiv. à une petite cat. et} \unsure{grd} devant $\pi$, $2^\circ)$ les $\varinjlim$ de type $I_{\alpha_0}$ existent dans $I$, $3^\circ)$ le foncteur $I\longrightarrow M$, $i\mapsto x_i$, commute aux $\varinjlim$ de type $I_{\alpha_0}$, et enfin $4^\circ)$ les $x_i\in\mathrm{Ob}\,M_\pi$, (comparer 4.13.11), on a ceci (où on désigne par $I_{\underline S}$ la sous-catégorie str. pleine de $I$ formée des $i$ tels que $x_i\in\mathrm{Ob}\,\underline S$) : on a
\[
 x\in\mathrm{Ob}\,\underline S\quad\text{\unsure{sss|ssi}}\quad I_{\underline S}\text{ est cofinal }\unsure{\text{dans}}\,I.
\]

c') \inserted{(\unsure{et|{} }\uuline{si} $I_{\alpha_0}$ est grd devant $\pi_0$ et $\pi^{\pi_0}=\pi$, pour pouvoir invoquer 4.13.14), cf. p. 355\unsure{$>$|.}} Comme c), \deleted{mais} en supposant $I$ \unsure{ordonnée|ordonné}.

% slanted
\margin{\deleted{\unsure{Attention|Ajouter|NB}}}

\margin{\uline{NB} Si on ne suppose pas $I_{\alpha_0}$ grd devant $\pi_0$ et $\pi^{\pi_0}=\pi$, alors seules les conditions de a) à c) (a, a', b, c) sont équivalentes, et \uline{impliquent} c', mais ne \uline{\unsure{lui|leur}} sont \unsure{peut être|peut-être} pas équivalentes.}

Démonstration. Elle se fait suivant le $=$ diagramme logique que celle de 4.13.15
\diagram{\begin{tikzcd}[ampersand replacement=\&,column sep=small,row sep=small]
 \& \& \& c'\arrow[dr,Rightarrow,"!"{name=bang}] \& \& {\begin{gathered}\text{ici, il faut supposer}\\ I_{\alpha_0}\text{ grd devant }\pi_0\\ \text{et }\pi^{\pi_0}=\pi\end{gathered}}\\
 a\arrow[r,Rightarrow] \& a'\arrow[r,Rightarrow,"!"] \& c\arrow[ur,Rightarrow]\arrow[dr,Rightarrow] \& \& a'\arrow[r,Rightarrow] \& {a,}\\
 \& \& \& b\arrow[ur,Rightarrow] \& \& {}
 \arrow[from=bang,to=1-6,end anchor=west]
\end{tikzcd}}
où on a marqué du signe ! les deux implications qui sont ``vaches'', en ce sens qu'elles font appel au lemme-clef 4.13.17 (p\unsure{g|ages} 329-330), qui se prouve par une construction transfinie un peu \unsure{vissée|vicieuse}. La démonstration \deleted{\unsure{suit}} des sept implications \inserted{du diagramme} est essentiellement celle donnée dans les pages précédentes 311-337,

% archive sheet number: 426

\scanpage{428}
\mspage{340}

à l'exception de $a'\Rightarrow a$ qui est un peu plus simple dans le cas actuel. On va la dégager sous forme de lemme un peu plus \unsure{général|générale} : (NB la dém. de $a'\Rightarrow a$ dans 4.13.15 avait été omise comme ``triviale'' ! cf. \unsure{pg|p.} 337 \unsure{\ldots})

\uline{Lemme} 4.13.19 : Soient $\unsure{M|\mathcal M}$ une catégorie accessible, $\underline S$ une sous-catégorie accessible de $M$, $\pi$ un cardinal utile pour $M$ et tel que $\underline S$ soit $\pi$-accessible, i.e. $\underline S$ est l'image essentielle de $\mathrm{Ind}_\pi(\underline S(\pi))$ par le foncteur pl. fidèle canonique $\mathrm{Ind}_\pi(\underline S(\pi))\longrightarrow M$, où $\underline S(\pi)=\underline S\cap M_\pi$.

On suppose qu'il existe un ordinal limite $\alpha$, tel que $\mathrm{card}\,\alpha\leq\pi$ et tel que $M_\pi$ et $\underline S\deleted{\unsure{{}_{\pi}|(\pi)}}(\pi)$ $(=S_\pi)$ soient stables par $\varinjlim$ de type $I_\alpha$, \deleted{Alors} et que le foncteur d'inclusion $\underline S(\pi)\longrightarrow M\deleted{\unsure{(\pi)}}_\pi$ y commute. Soit $\unsure{\Phi|\mathscr I|\mathscr D}$ un ensemble de catégories \deleted{de \unsure{cardinaux|cardinal}} telles que $\forall I\in\Phi$, on ait $\unsure{\mathrm{card}\,\mathrm{red}\,I|\mathrm{card}\,\mathrm{rel}\,I|\mathrm{card}\,\text{réel}\,I}\leq\pi$, et supposons que $M\deleted{\unsure{(\pi)}}_\pi$ \deleted{soit} \deleted{et $\underline S(\pi)$ soient stables} soit stable par $\varinjlim$ de type $\Phi$, et que $\underline S(\pi)$ soit stable dans $M_\pi$ par lesdites limites (ou ce qui revient au $=$, qu'il soit stable en lui-même pour ces $\varinjlim$\unsure{\textsuperscript{t}|,}, et que le foncteur d'inclusion y commute). Alors $M$ et $\underline S$ sont stables par $\varinjlim$ de type $\Phi$, et le foncteur d'inclusion $\underline S\longrightarrow M$ y commute, i.e. $M$ est stable pour ces limites, et $\underline S$ stable dans $M$ pour ces \unsure{$\Rightarrow$|$=$} limites.

% archive sheet number: 427

\scanpage{429}
\mspage{341}

\uline{Démonstration} On applique 4.11.1 \inserted{(p. 267-268)} \unsure{séparément|séparement} à $\unsure{M|\mathcal M}\approx\mathrm{Ind}_\pi M\deleted{(\pi)}_\pi$ et à $S\approx\mathrm{Ind}_\pi\,\underline S(\pi)\unsure{\}|)}$, pour conclure que $M$ et $\underline S$ sont stables séparément par $\varinjlim$ de type $\unsure{\Phi|\mathscr I|\mathscr D}$. Il reste à voir que le foncteur $\underline S\longrightarrow M$ y commute. Cela se voit comme dans la démonstration de 4.11.1, \unsure{on} obtient plus généralement ceci :

\uline{Lemme} 4.13.20 (cor. à 4.11.1) : \unsure{soient|soient\textsuperscript{t}} $f_0:C\longrightarrow C'$, un foncteur entre catégories \unsure{équivalentes à|équiv. à} des petites catégories. $\pi$ un cardinal infini, $f:M\overset{\text{déf}}{=}\mathrm{Ind}_\pi C\longrightarrow M'\overset{\text{déf}}{=}\mathrm{Ind}_\pi C'$
% slanted
\margin{Il vaut mieux \unsure{prendre|parler de} un foncteur $f_0:C\longrightarrow N$, et dire que $f_0$ commutant aux $\varinjlim$ de type $\Phi$, \unsure{alors} il en est de $=$ de son \unsure{extension|extens.} $\mathrm{Ind}_\pi(C,N)\longrightarrow N$. (NB on suppose $N$ satisfait $\mathrm L_\pi$).}
l'extension canonique $\mathrm{Ind}_\pi(f_0)$ de $f_0$. Soit $\Phi$ \inserted{et $\alpha$} comme dans 4.12.1, et supposons que $\deleted{\unsure{I_\alpha}}I_\alpha\in\Phi$ et que $f_0$ commute aux $\varinjlim$ de type $\Phi$. Alors $f$ commute à ces $\unsure{\Rightarrow|=}$ limites inductives.

\uline{Corollaire} 4.13.21. \inserted{Dans les conditions de 4.13.20,} Supposons que $\Phi$ contienne les catégories de cardinal $\leq\pi$ (resp. les catégories filtrantes de cardinal $\leq\pi$, \inserted{resp. des cat. filtr. \unsure{grdes|grds} devant $\pi_0$ de card. $\leq\pi$,} resp. les catégories discrètes de cardinal $\leq\pi$). Alors $f$ commute aux \inserted{petites} $\varinjlim$ \deleted{quelconques} (resp. aux petites $\varinjlim$ filtrantes, \inserted{resp. aux petites $\varinjlim$ filtrantes grandes devant $\pi_0$,} resp. aux petites sommes).

(Réciproque claire et ``laissée au lecteur''.)

Se voit comme 4.4.14 et 4.4.14.1 (p. 97,98).

% archive sheet number: 428

\scanpage{430}
\mspage{342}

\deleted{\uline{Définition} 4.13.22. Soient $\unsure{M|\mathcal M}$ une catégorie accessible, $\underline S$ une sous-catégorie de $M$, $\pi$ un cardinal \deleted{infini, utile} infini, $\unsure{\Phi|\mathscr I|\mathscr D}$ un ens. de \unsure{catégorie|catégories}. On suppose que \deleted{\unsure{on ait}} pour tt $I\in\Phi$, on ait $\unsure{\mathrm{card\,red}\,I|\mathrm{card\,rel}\,I}\leq\pi$, \deleted{On dit que $\underline S$ est $\Phi$-$\pi$-accessible dans $M$, si} et qu'il existe un ordinal limite $\alpha_0$ tel que $I_{\alpha_0}\in\Phi$.}

Voici enfin, il me semble, la forme définitive (?) de 4.13.15, 4.13.18 :

% slanted
\margin{cf Variante 4.13.3\unsure{6|5} p 358}

\uline{Théorème} 4.13.22. \deleted{\unsure{Soit}}Soient $M$ une catégorie accessible, \inserted{$\pi$ un cardinal infini, utile pour $M$,} et \deleted{\unsure{$S$}}$\Phi$ un ens. de catégories satisfaisant les deux conditions

\unsure{$1^\circ)$|a)} $\deleted{I}J\in\Phi$ implique $\mathrm{card}\,\unsure{\mathrm{rd}|\mathrm{red}}\,\deleted{I}J\leq\pi$

\unsure{$2^\circ)$|b)} $\exists$ ordinal \uline{limite} $\alpha_0$ avec $I_{\alpha_0}\in\Phi$.

\inserted{On suppose $M$ stable par $\varinjlim$ de \deleted{\unsure{du type}} $\Phi$.} Soit $\underline S$ une sous-catégorie \uline{str. pleine} de $M$, et $\underline S(\pi)=\underline S\cap M_\pi$. \deleted{\unsure{Consid}} Désignons par $\mathrm L(\Phi)$ la condition sur $\underline S(\pi)$ : $\underline S(\pi)$ est stable dans $\unsure{M_\pi|M}$ (ou ce qui revient au $=$, dans $M_\pi$) par $\varinjlim$ de type $\Phi$. Ceci posé, les conditions suivantes \inserted{sur $\underline S$} sont équivalentes :

\textcircled{a} $\underline S$ est $\pi$-accessible dans $M$, et stable \deleted{\unsure{par|pour}}\unsure{dans} $M$ par $\varinjlim$ de type $\Phi$.

a') $\underline S$ est $\pi$-accessible dans $M$, et \inserted{$\underline S(\pi)$} satisfait à $\mathrm L(\Phi)$.

% archive sheet number: 429

\scanpage{431}
\mspage{343}

b) $\underline S(\pi)$ satisfait à $\mathrm L(\unsure{\Phi|\mathscr I|\mathscr D})$, et pour tt $x\in\mathrm{Ob}\,\unsure{M|\mathcal M}$, on a ceci : $x\in\mathrm{Ob}\,\underline S$ \unsure{sss|ssi} la sous-catégorie str. pleine $\underline S(\pi)_{/x}$ de $M(\pi)_{/x}$ est cofinale dans $M(\pi)_{/x}$.

\textcircled{c} $\underline S(\pi)$ satisfait à $\mathrm L(\Phi)$, et pour tt $x\in\mathrm{Ob}\,M$, et \deleted{de} \deleted{\unsure{t|tt}} \unsure{représentation} $x\simeq\varinjlim_{i\in I}x_i$ qui est $\Phi$-$\pi$-adaptée (cf. déf. plus bas), on a ceci : $x\in\mathrm{Ob}\,\underline S$ sss $I_{\underline S}$ est cofinal dans $I$, où $I_{\underline S}$ est la sous-catégorie str. pleine de $I$ formée des $i\in\mathrm{Ob}\,I$ tels que $x_i\in\mathrm{Ob}\,\deleted{\unsure{I|S}}\unsure{\underline S|\underline I}$. (ou encore, $x_i\in\mathrm{Ob}\,\underline S(\pi)$).

NB On dit qu'un système inductif $(x_i)_{i\in I}$ de $M$ est $\Phi$-$\pi$-adapté si 1) $I$ équiv. à une petite catégorie et $I$ \unsure{grde|grd} devant $\pi$ 2) les $\varinjlim$ de type $\Phi$ existent dans $I$, 3) le foncteur $i\unsure{\mapsto|-}x_i$ de $I$ dans $M$ commute aux $\varinjlim$ de type $\Phi$, et enfin 4) les $x_i\in\mathrm{Ob}\,M_\pi$.

c') ((\deleted{Si} $\pi^{\pi_0}=\pi$, où $\pi_0$ est un cardinal $\leq\unsure{\pi|\mathrm r}$ utile \unsure{à|pour} $M$, \deleted{\unsure{si on a} $\unsure{I_{\alpha_0}}$ de plus $I_{\alpha_0}$ est \unsure{grde} devant $\pi_0$, et enfin} et si \deleted{$\unsure{\Phi\ni I_{\alpha_0}}$ est} \deleted{\unsure{est}} $J\in\Phi$ implique que $J$ est grd devant $\pi_0$) Comme c), avec $I$ \uline{ordonnée}.

% slanted
\margin{\uline{NB} Quand les conditions restrictives ci-contre ne sont pas satisfaites, alors a, a', b, c sont équiv. et impliquent c', mais ne sont peut-être pas \unsure{équivalentes} à c').}

\margin{\uline{NB} Il faudrait \unsure{réénoncer|réén.} \unsure{p.|pl.} la condition en \unsure{gr|gros|$\gamma$} (si \unsure{\ldots}), cf p. 359\unsure{,|-}360}

\margin{NB Il y aurait intérêt à remonter cette définition (\uline{sans} b), déjà avant 4.13.15.}

\uline{Définition} 4.13.23. Soient $M$ une catégorie accessible, $\pi$ \inserted{\unsure{est de}} un cardinal, $\Phi$ un ens. de catégories, satisfaisant les deux conditions

\unsure{a*)|a)} $\deleted{\unsure{J|I}}\unsure{I|J}\in\Phi\Longrightarrow\mathrm{card}\,\unsure{\mathrm{red}|\mathrm{rel}}\,J\leq\pi$

\deleted{\unsure{b)}}b) $\exists$ ordinal limite $\alpha_0$, tel que $I_{\alpha_0}\in\Phi$.

% slanted
\margin{Si c'est pour poser une \uline{définition}, on peut se \unsure{contenter|contenter de} \unsure{a)|de a)} \unsure{sans|pour} \unsure{b)}. $\leftarrow$ $\nwarrow$}

% archive sheet number: 430

\scanpage{432}
\mspage{344}

Soit $S$ une partie de $\mathrm{Ob}\,\unsure{M|\mathcal M}$, $\underline S$ la sous-catégorie pleine qu'elle définit. On dit que $S$ est une \uline{partie} $\unsure{\Phi|\mathscr I|\mathscr D}$-$\pi$-\uline{accessible} de $\mathrm{Ob}\,M$, ou que $\underline S$ est une \uline{sous-catégorie} $\Phi$-$\pi$-\uline{accessible} \uline{dans} $M$, si \unsure{,} les conditions suivantes sont satisfaites

\margin{\unsure{condition|conditions} prél. sur $M$ $\leftarrow$}

$1^\circ)$ $\pi$ est utile pour $M$, $M$ est stable par $\varinjlim$ de type $\Phi$.

$2^\circ)$ $\underline S$ est une sous-catégorie $\pi$-accessible \unsure{:} dans $M$.

$3^\circ)$ $\underline S$ est \deleted{accessible} stable dans $M$ par $\varinjlim$ de type $\Phi$.

% slanted
\margin{\uline{NB} Quand $\Phi=\{I_{\alpha_0}\}$, on dit aussi que $\underline S$ est $\alpha_0$-$\pi$-accessible dans $M$. Si $\Phi$ est formé des cat. de cardinal $\leq\pi$ et grdes devant $\pi_0$, on dit que $\underline S$ est $\pi_0$-$\pi$-\unsure{accessible} dans $M$.}

\uline{NB} Quand la condition préliminaire $1^\circ)$ \uline{sur $M$} est satisfaite, \inserted{la validité} des conditions $2^\circ)$ $3^\circ)$ \unsure{équivaut|équivalent} donc à la condition que $\underline S$ soit \uline{strictement} pleine dans $M$, et que \inserted{de plus} $\underline S$ satisfasse la condition a) du th. 4.13.22. D'autre part, \unsure{cette|ces} condition est équivalente à \deleted{aux} \inserted{chacune des} \unsure{conditions} a', b, c de ce même théorème, et aussi à la dernière condition c'), \unsure{sous} réserve qu'on ait $\pi^{\pi_0}=\pi$ (pour un cardinal $\pi_0<\pi$ utile pour $M$) et que $J\in\Phi$ implique que $J$ est \unsure{grde|grd} devant $\pi_0$.

Conceptuellement, c'est la définition donnée dans 4.13.23, \deleted{qui \unsure{p}} i.e. la condition \textcircled{a} de 4.13.22, qui paraît \unsure{le|la} plus \unsure{clair|claire}, ou sa variante immédiate a'\deleted{\unsure{a)|)}}. Mais du point de vue pratique, du \deleted{\unsure{le mani}} \unsure{travail} avec les parties accessibles, c'est la propriété \textcircled{c} qui paraît le plus commode. Le poids principal du th. 4.13.22 est bien

% archive sheet number: 431

\scanpage{433}
\mspage{345}

l'implication $a\Longrightarrow c$ (via $a\Rightarrow a'\overset{!}{\Rightarrow}c$). Donc Il faut surtout voir ce théorème comme donnant un critère remarquable d'appartenance à $\mathrm{Ob}\,\underline S$ d'un objet $x$ de $\unsure{M|\mathcal M}$ qui est présenté par un syst. inductif $\unsure{\Phi|\mathscr I|\mathscr D}$-$\pi$-adapté :
\[
x\unsure{+|\in}\mathrm{Ob}\,\underline S\quad\Longleftrightarrow\quad I_{\underline S}\text{ cofinal},\tag*{\unsure{$(*)$|$(**)$}}
\]
et dans cette équivalence, c'est l'implication $\Longrightarrow$ qui est la partie \deleted{``vache''} \inserted{non-triviale}, et remarquable.

\uline{Proposition} 4.13.\deleted{\unsure{23}}24. Soient \deleted{\unsure{$\Phi$}}$M$, $\pi$, $\Phi$ comme dans 4.13.23, \unsure{$*$} on suppose que \deleted{$M$ \unsure{stable}} la condition \uline{1°} de \deleted{\unsure{la dé}} 4.13.23 est satisfaite ($\pi$ utile pour $M$, $M$ stable par $\varinjlim$ de type $\Phi$).

% slanted
\margin{NB. La condition $I_{\alpha_0}\in\Phi$ n'est pas utilisée ici.}

Soit $\Sigma_\pi^{\prime\Phi}\deleted{=\unsure{\Sigma(M,\pi,\Phi)}}$ l'ens des sous-catégories \deleted{que} \inserted{$\underline S$} de $M$ qui sont \deleted{\unsure{$\underline S$}} $\Phi$-$\pi$-accessibles, dans $M$, et $\Sigma^{\Phi}_\pi$ l'ens des sous-catégories \inserted{$\underline S(\pi)$} de $M_\pi$ qui sont str. pleines, \deleted{\unsure{Karou}} et stables dans $M_\pi$ par $\varinjlim$ de type $\Phi$. (Ces sous-catégories sont forcément Karoubiennes, car \deleted{$I_{\alpha_0}$} stables par $\varinjlim$ de type $I_{\alpha_0}$, donc aussi par facteurs directs). Alors les bijections (4.13.\deleted{\unsure{1}}4.1) (p. \deleted{\unsure{p}} 288) induisent des bijections inverses l'une de l'autre
\diagram{\begin{tikzcd}[ampersand replacement=\&]
 \Sigma_{\pi}^{\Phi}\arrow[r,shift left] \& {\Sigma_{\pi}^{\prime\Phi}.}\arrow[l,shift left]
\end{tikzcd}}

% archive sheet number: 432

\scanpage{434}
\mspage{346}

\uline{Proposition} 4.13.25. Soient $\unsure{M|\mathcal M}$, $\pi$, $\unsure{\Phi|\mathscr I|\mathscr D}$ comme dans 4.13.23, avec $I_{\alpha_\pi}\in\Phi$ ($\alpha_\pi$ étant le plus petit ordinal \unsure{de|du} cardinal $\unsure{\Phi|\pi}$). Alors l'ens $\Sigma_\pi^{\prime\Phi}$ des sous-\unsure{catégorie|catégories} de $M$ $\Phi$-$\pi$-accessibles dans $M$ est stable par intersection de familles de cardinal $\leq\pi$.$^{*}$

% slanted
\margin{\unsure{Si|si} $\pi$ n'est pas un cardinal limite, il faut dire : familles de cardinal $\deleted{\unsure{\leq|{}=}}\!<\pi$. cf. 4.13.10.3., où se \unsure{trouve} un énoncé relatif, légèrement plus général\unsure{;|.}}

\deleted{\unsure{Les}} (\inserted{NB.} \unsure{la|les} bijection $\Sigma_\pi^\Phi\simeq\Sigma_\pi^{\prime\Phi}$ de 4.13.24 \unsure{respecte|respectent} les relations d'ordre naturelles d'inclusion, donc transforme intersections en intersections)

Se voit comme 4.13.7 (qui est contenu dans l'énoncé actuel), en utilisant le lemme 4.13.10.

\uline{4.13.26} Soient $M$, $\pi$, $\Phi$ comme dans 4.13.23, et soit $\underline S$ une sous-catégorie \uline{str. pleine} de $M$, $f:\underline S\to M$ le foncteur d'inclusion. \deleted{On suppose déjà que $M$ satisfait la \unsure{1°|1\textsuperscript{re}} de la \unsure{partie|page} \unsure{qui|préc.} \unsure{loc. cit.} \unsure{i.e.|$\pi$}, $M$ \unsure{est|et}} La condition que $\underline S$ soit $\Phi$-$\pi$-accessible dans $M$ se décompose en :

1) $\pi$ est utile pour $M$ \uline{et} pour $\underline S$.

2) $M$ et $\underline S$ sont stables par $\varinjlim$ de type $\Phi$.

3) $f$ commute aux dites limites.

4) $f$ est $\pi$-accessible

5) $f(\unsure{S|\underline S}_\pi)\subset M_\pi$.

% archive sheet number: 433

\scanpage{435}
\mspage{347}

Soit donc $\pi'>\pi$ un autre cardinal, quand peut-on assurer que \inserted{si} $\underline S$ est $\unsure{\Phi|\mathscr I|\mathscr D}$-$\pi$-accessible dans $\unsure{M|\mathcal M}$, il soit aussi $\Phi$-$\pi'$-accessible ?

Les conditions 2) 3) 4) sont visiblement stables par passage de $\pi$ à $\pi'$. Restent les conditions 1) et 5). On sait que si ${\pi'}^{\pi}=\pi'$, alors la condition 1) reste valable pour $\pi'$. Plus \unsure{généralement|généralemt}, si $\pi_0^{\deleted{\unsure{\prime|*}}}$ est un cardinal $\leq\pi$ tel que \unsure{lui|le} soit utile pour $M$ et pour $\underline S$, alors si $\pi'>\pi_0$ est tel que ${\pi'}^{\pi_0}=\pi'$, alors $\pi'$ est aussi utile pour $M$ et pour $\underline S$.

Si on admet l'\unsure{HCC|HCG} (hypothèse du continu généralisée), il en est ainsi pour tt $\pi'>\pi_0$, donc pour \unsure{tt|tout} $\pi'>\pi$.

Reste la condition 5), pour laquelle une réponse est donnée par 4.12.1 (p. 27\deleted{\unsure{2|3}}\unsure{3|2}). Si $\pi_0$ \inserted{$\leq\pi$} est choisi de façon que $f$ soit $\pi_0$-accessible (\unsure{p.e.|p.ex.} $\pi_0=\pi$), et si \deleted{$\pi_1\geq\pi$ $(\unsure{>|\geq}\pi_0)$ \unsure{on pose} est tel que $\pi_1=\sup_{\unsure{a}} f(\unsure{\in|{}<}M_{\pi_0})$ \unsure{ou} $\pi_1=\mathrm{Card}\,\unsure{\mathrm{red}|\mathrm{rel}|\mathrm{réel}}\,\underline S_\pi$,}
\[
\pi'_1=\mathrm{Card}\,\unsure{\mathrm{red}|\mathrm{rel}|\mathrm{réel}}\,\underline S_{\pi_0}\quad(\leq\mathrm{Card}\,\mathrm{red}\,\underline S(\pi)),
\]
\deleted{on voit que}
\[
\pi''_1=\mathrm{Sup}\,(\pi'_1,\pi),
\]
on voit que si $\pi'\geq\pi''_1$ (i.e. $\pi'\geq\pi'_1,\pi$) est tel que $\deleted{\unsure{\pi'\geq\lambda|\pi>\lambda|\pi'}}\ {\pi'}^{\pi_0}=\pi'$, alors on a bien
\[
 f(\underline S_{\pi'})\subset M_{\pi'}.
\]

% archive sheet number: 434

\scanpage{436}
\mspage{348}

Si donc \inserted{de plus} $\pi'>\pi\inserted{\,\geq\pi_0}$ (comme supposé au début), et si on suppose l'\deleted{\unsure{hyp.}}\unsure{HCG|HCC}, \deleted{\unsure{donc}} \unsure{donc|alors} la condition ${\pi'}^{\pi_0}=\pi'$ est automatique, et il reste la seule condition $\pi'\geq\pi''_1$ i.e. en définitive $\pi'\geq\mathrm{Card}\,\unsure{\mathrm{red}|\mathrm{rel}|\mathrm{réel}}\,\underline S_{\pi_0}$. Pour résumer :

\uline{Proposition} 4.13.27. Sous les conditions de 4.13.23 (\unsure{sans|sous} condition b) sur $\unsure{\Phi|\mathscr I|\mathscr D}$), soit $\underline S$ une sous-catégorie de $\unsure{M|\mathcal M}$ qui est $\Phi$-$\pi$-\unsure{adaptée|adapté}. Soit $\pi_0\leq\pi$ un cardinal \deleted{tel que} qui est utile pour $\underline S$ et pour $\underline M$, \deleted{\unsure{que ce}} et tel que $f:\underline C\to\underline M$ soit \unsure{$\pi_0$-accessible|$\pi_0$-accessibl} i.e. tel que $\underline S$ soit stable dans $M$ par $\varinjlim$ \unsure{grdes|grandes} devant $\pi_0$ (on peut p.ex. prendre $\pi_0=\pi$). Alors pour tout cardinal $\pi'>\pi$ tel que a) ${\pi'}^{\pi_0}=\pi'$ et b) $\pi'\geq\mathrm{Card}\,\mathrm{red}\,\underline S_{\pi_0}$), \deleted{\unsure{alors|on a}} $\underline S$ est $\Phi$-$\pi'$-accessible dans $M$.

Donc si $\pi'>\pi$ est tel que ${\pi'}^{\pi}=\pi'$, $\pi'>\mathrm{Card}\,\mathrm{red}\,\underline S_\pi$, alors $\underline S$ est $\Phi$-$\pi'$-\unsure{accessible|accessibl}

(Il suffit p.ex. de prendre $\pi'=2^c$, avec $c\geq\mathrm{Sup}\,(\pi,\mathrm{Card}\,\mathrm{red}\,\underline S(\pi))$, ou $\pi'=2^c$ avec $c\geq\pi$ et $2^c\geq\mathrm{card}\,\mathrm{red}\,\underline S(\pi)$.)

Si on admet l'\unsure{HCC|HCG}, il suffit donc qu'on ait $\pi'\geq\mathrm{Sup}\,(\pi,\mathrm{Card}\,\mathrm{red}\,\underline S(\pi))$.

% archive sheet number: 435

